% Source-derived chapter 08: Compactified parameter spaces
% Original source: rigid-local-systems-multiplicative-eigenvalue-problem
\chapter{Compactified parameter spaces}
\label{rigid-local-systems-multiplicative-eigenvalue-problem-chapter-08}
\source{rigid-local-systems-multiplicative-eigenvalue-problem}

\begin{remark}[Source section]
\label{rigid-local-systems-multiplicative-eigenvalue-problem-chapter-08-source}
\source{rigid-local-systems-multiplicative-eigenvalue-problem}
\source{rigid-local-systems-multiplicative-eigenvalue-problem}
This chapter reproduces the corresponding section of the retrieved
LaTeX source, including its definitions, statements, examples, and
proofs. It has not yet been translated into Lean.
\notready
\end{remark}

% The source section title was: Compactified parameter spaces
\label{compact}
\source{rigid-local-systems-multiplicative-eigenvalue-problem}
The remaining aims for this paper are the following: completing the proofs of Theorems \ref{Adagio},\ref{brahms}, and Proposition \ref{condensed}, and  setting up the induction apparatus to induce extremal rays from Levi subgroups giving rise to the maps \eqref{productstructure} and \eqref{inductiones}.

Assume that the generalized Gromov-Witten number $\langle\sigma_{I^1}, \sigma_{I^2},\dots,\sigma_{I^s}\rangle_{d,D}=1$, where $D=-\deg(\mn)$ (Definition  \ref{GWgen} (We use this general set up with $D$ possibly non-zero because we will need to use shift operations to handle the case $1\in I^j$ for some $j$). This implies that $\Omega(d,r,\mn,n,\vec{I})\to \Par_{n,\mn,S}$ is birational and
\begin{equation}\label{timer1}
\source{rigid-local-systems-multiplicative-eigenvalue-problem}
dn +\deg(\mn)r +r(n-r)-\sum_{i=1}^s |\sigma_{I^i}|=0.
\end{equation}

Recall that  $\Omega(d,r,\mn,n,\vec{I})$ is a compactification over $\Par_{n,\mn,S}$ of $\Omega^0(d,r,\mn,n,\vec{I})$ (which is the set of subbundles of the universal bundle of determinant $\mn$ and rank $n$ with fibers in given Schubert states). $\Omega(d,r,\mn,n,\vec{I})$ has several defects: (1) for a  point $(\mv,\mw,\me,\gamma)$ not in $\Omega^0(d,r,\mn,n,\vec{I})$,  $\mw/\mv$ can have torsion and one does not have an inductive structure of a  subbundle and quotient which are both bundles with induced flags, (2) even when $\mw/\mv$ is locally free the induced flag structures cannot be assigned without discontinuities ,and (3) $\Omega(d,r,\mn,n,\vec{I})$ could be singular.

For the induction map and the proof of Proposition \ref{fulC} we need a relative
partial compactification of $\Omega^0(d,r,\mn,n,\vec{I})\to \Par_{n,\mn,S}$ without these defects, as well as stacks for the relevant Levi subgroups. We will also need to understand the ramification properties of $\Omega(d,r,\mn,n,\vec{I})\to \Par_{n,\mn,S}$.  This will result in a (known) multiplicative generalization of a conjecture of Fulton, which we prove for completeness here (see Proposition \ref{fulC} below).



\subsection{Stacks for Levi subgroups}
\begin{defi}
\source{rigid-local-systems-multiplicative-eigenvalue-problem}
\notready\label{parul}
\source{rigid-local-systems-multiplicative-eigenvalue-problem}
Let $\ParL=\ParL(-d,r,n-r,d-D)$ be the moduli of tuples $(\mv,\mq,\mf,\mg,\gamma)$ where
\begin{enumerate}
\item $\mv$ (resp. $\mq$) is a degree $-d$ (resp. $d-D$), rank $r$ (resp. rank $n-r$) vector bundle on $\Bbb{P}^1$,
\item  $\mf\in \Fl_S(\mv)$, $\mg\in \Fl_S(\mq)$,
\item $\gamma$ is an isomorphism $\det \mv\otimes \det\mq\leto{\sim}\mn=\mathcal{O}(-D)$.
\end{enumerate}
\end{defi}
There is a diagram
\begin{equation}\label{Dbasic}
\source{rigid-local-systems-multiplicative-eigenvalue-problem}
\xymatrix{
 & \Omega^0(d,r,\mn,n,\vec{I})\ar[dl]^{\pi}\ar[dr]\\
\Par_{n,\mn,S}   &  & \ParL\ar@/_/[ul]_i\ar[ll]^{i'=\pi\circ i}
}
\end{equation}
Here $i$ is a ``direct sum morphism", $i(\mv,\mq,\mf,\mg,\gamma)= (\mw,\me,\gamma')\in \Par_{n,\mn,S}$ where $\mw=\mv\oplus\mq$, $\gamma': \det(\mw)=\det(\mv)\otimes\det(\mq)\leto{\gamma} \mathcal{O}(-D)$, and $\me$ is given by the following rule. For $p=p_i$, write $I^i=\{i_1>i_2>\dots>i_r\}$. Define $E^p_{\bull}$ as follows: For  $1\leq k\leq n$, define
$$E^p_{j}=F^p_{a}\oplus G^p_{j-a}$$
where  $i_a\leq j <i_{a+1}$ (with $i_0=0$ and $i_{r+1}=n$). Let $\mathcal{R}\subseteq  \Omega^0(d,r,\mn,n,\vec{I})$ be the ramification locus of $\pi$ (see Section \ref{RamDiv}).

\subsection{Desired properties of partial compactifications}
One of our basic techniques, inspired by Ressayre's work \cite[Section 4.1]{R1} (see \cite[Section 8.6]{BHermit} for a comparison)  is to take line bundles on $\ParL$, and lift them to $\Omega^0(d,r,\mn,n,\vec{I})$.  By Zariski's main theorem, $\Omega^0(d,r,\mn,n,\vec{I})-\mathcal{R}$ is an open substack of $\Par_{n,\mn,S} $, and we thus get a line bundle on an open substack of $\Par_{n,\mn,S} $.

\begin{remark}
\source{rigid-local-systems-multiplicative-eigenvalue-problem}
\notready\label{explainmore}
\source{rigid-local-systems-multiplicative-eigenvalue-problem}
This open substack misses all the basic divisors $D(a,j)$ from the introduction (in the case $\mn=\mathcal{O}$).
\end{remark}
By the above remark the complement of this open subset may have codimension 1 components, and we cannot then naturally extend the line bundle to  $\Par_{n,\mn,S}$. One looks for an enlargement of   $\Omega^0=\Omega^0(d,r,\mn,n,\vec{I})$ that does not have this defect (and which still admits maps to $\Par_{n,\mn,S} $ and $\ParL$).

In fact this happens in the simplest possible fashion whenever $1\not\in I^i$ for any $i$, and the shift operations allow us to reduce the general case to this case. This partial compactification will be inside $\Omega=\Omega(d,r,\mn,n,\vec{I})$.

 We do this in two steps: First we enlarge to $\widehat{\Omega}^0\to \Par_{n,\mn,S}$ where we only weaken the Schubert conditions, the subsheaves remaining subbundles. The map   $\Omega^0\to \ParL$ will be shown to extend to $\widehat{\Omega}^0\to \ParL$ when $1\not\in I^i$ for any $i$.

The second step will extend this to a substack  $\widehat{\Omega}_{S}^0\subseteq \Omega$ over $\Par_{n,\mn,S}$, where the subbundle condition is imposed only at points of $S$. We will then have
$$\Omega^0\subseteq \widehat{\Omega}^0\subseteq \widehat{\Omega}_{S}^0\subseteq \Omega.$$ The map $\widehat{\Omega}^0\to \ParL$ will not be shown to extend to $\widehat{\Omega}_{S}^0\to \ParL$. Instead we will show in Proposition \ref{bee} that   $\widehat{\Omega}_{S}^0-  \widehat{\Omega}^0$ lies inside the ramification divisor $\mathcal{R}$, but  irreducible components of $\Omega-\widehat{\Omega}_{S}^0$ are of codimension $\geq 2$. The desired enlargement of $\Omega^0(d,r,\mn,n,\vec{I})$ will then be $\widehat{\Omega}^0$ (and not $\widehat{\Omega}_{S}^0$!).
\subsection{Construction of the partial compactification}
For convenience, we weaken our assumptions to the following:
Suppose $\vec{I}=(I^1,\dots,I^s)$ is an $s$-tuple of subsets  of $[n]$ each of cardinality $r$. Let $\Omega=\Omega(d,r,\mn,n,\vec{I})$.  The representable morphism $\Omega\to \Par_{n,\mn,S}$ has relative dimension $dn +\deg(\mn)r +r(n-r)-\sum_{i=1}^s |\sigma_{I^i}|$. We will make the assumption that this relative dimension is zero:
\begin{equation}\label{timer}
\source{rigid-local-systems-multiplicative-eigenvalue-problem}
dn +\deg(\mn)r +r(n-r)-\sum_{i=1}^s |\sigma_{I^i}|=0.
\end{equation}

\subsubsection{Definition of $\widehat{\Omega}^0=\widehat{\Omega}^0(d,r,\mn,n,\vec{I})$}\label{bruckner92}
\source{rigid-local-systems-multiplicative-eigenvalue-problem}
We create a  smooth open subset of $\Omega(d,r,\mn,n,\vec{I})$ containing $\Omega^0(d,r,\mn,n,\vec{I})$ in the following way.

\begin{defi}
\source{rigid-local-systems-multiplicative-eigenvalue-problem}
\notready\label{ToI}
\source{rigid-local-systems-multiplicative-eigenvalue-problem}
\hspace{2em}
\begin{enumerate}
\item For a subset $I$ of $[n]$ of cardinality $r$, let $T(I)\subset [n]$ be the subset of all $j\in [n]$ such that one of the following holds $j\in I$ or $j+1\not\in I$ (equivalently $a\not\in T(I)$ if and only if $a\not\in I$ and $a+1\in I$). Note that $T(I)$ contains $n$.
\item In the setting of Definition \ref{ineffable}, define a smooth variety $\widehat{\Omega}^0_I(E_{\bull})$, with $\Omega^0_I(E_{\bull})\subseteq \widehat{\Omega}^0_I(E_{\bull})\subseteq {\Omega}_I(E_{\bull})$ by
$$\widehat{\Omega}^0_I(E_{\bull})=\{V\in \Gr(r,W)\mid \rk (V\cap E_j)=a, \text{ if } i_a\leq j<i_{a+1}\text { and } j\in T(I)\}.$$
It is easily seen that  $\widehat{\Omega}^0_I(E_{\bull})$ carries a transitive action of a connected group and is hence smooth. It is easy to see that all irreducible components of
${\Omega}_I(E_{\bull})-\widehat{\Omega}^0_I(E_{\bull})$ are of codimension $\geq 2$ in ${\Omega}_I(E_{\bull})$, see \cite[Lemma 7.9]{BHermit}.
\end{enumerate}
\end{defi}
Define $\widehat{\Omega}^0=\widehat{\Omega}^0(d,r,\mn,n,\vec{I})\subseteq \Omega(d,r,\mn,n,\vec{I})$ to be the smooth substack parametrizing tuples $(\mv,\mw,\me,\gamma)$ with $\mv$ a subbundle of $\mw$, and $\mv_p\in \widehat{\Omega}^0_{I^i}(\me_p)\subseteq \Gr(r,\mw_p)$ for all $p=p_i\in S$. Clearly, $\widehat{\Omega}^0(d,r,\mn,n,\vec{I})\supseteq {\Omega}^0(d,r,\mn,n,\vec{I})$.
By \cite[Lemma 7.10]{BHermit}), the morphism  $\Omega^0 \to \ParL$ extends to a morphism $\widehat{\Omega}^0\to\ParL$, and we get a commutative diagram
 \begin{equation}\label{morph1}
\source{rigid-local-systems-multiplicative-eigenvalue-problem}
\xymatrix{
 \Omega^0 \ar[r]\ar[d] &  \widehat{\Omega}^0\ar[dl]\\
\ParL
}
\end{equation}
\subsubsection{Definition of $\widehat{\Omega}_{S}^0=\widehat{\Omega}_{S}^0(d,r,\mn,n,\vec{I})$}

There is an even larger  open subset: Instead of requiring that $\mv$ is a subbundle of $\mw$, we can require only that it is a subbundle near points of $S$. Denote this open subset by $\widehat{\Omega}_{S}^0=\widehat{\Omega}_{S}^0(d,r,\mn,n,\vec{I})$ with the subscript $S$ denoting that the subbundle property is only imposed at points of $S$. Clearly $$\Omega^0\subseteq \widehat{\Omega}^0\subseteq \widehat{\Omega}_{S}^0\subseteq \Omega.$$
\begin{proposition}
\source{rigid-local-systems-multiplicative-eigenvalue-problem}
\notready\label{bee}
\source{rigid-local-systems-multiplicative-eigenvalue-problem}
\hspace{2em}
\begin{enumerate}
\item[(a)] If $1\not\in I^i$ for all $i$, then   all irreducible components of $\Omega-\widehat{\Omega}_{S}^0$ are of codimension $\geq 2$.
\item[(b)] $\widehat{\Omega}_{S}^0$ is a smooth Artin stack.
\item[(c)] Any point in $\widehat{\Omega}_{S}^0-\widehat{\Omega}^0$ is in the ramification locus of $\widehat{\Omega}_{S}^0\to \Par_{n,\mn,S}$.
%\item[(d)]If $1\not\in I^i$ for all $i$, then $\widehat{\Omega}^0-\mathcal{R}$ maps isomorphically to a open substack of $\Par_{n,\mn,S}$ whose complement has codimension $\geq 2$.
\end{enumerate}
\end{proposition}
Part (a) in the above theorem follows from the computations in Sections \ref{bruck02} and \ref{bruck03}.
Note that if $J$ is as in Section \ref{bruck03}, then ${\Omega}^0_J(E_{\bull}) \subseteq \widehat{\Omega}^0_{I^1}(E_{\bull})$, and therefore $\widehat{\Omega}_{S}^0(d,r,\mn,n,\vec{I})$ includes the codimension 1 degenerations of Section \ref{bruck03}.
 $\widehat{\Omega}_{S}^0$ is smooth and representable over $\op{Quot}(d,r,\mn,n)$ which implies (b).
 Part (c) of Proposition \ref{bee} is proved in  Section \ref{martel}.
\subsection{Ramification divisors}\label{RamDiv}
\source{rigid-local-systems-multiplicative-eigenvalue-problem}
Suppose $\mathcal{X}\leto{\pi} \mathcal{Y}$ is a representable morphism of smooth Artin stacks of the same dimension over $\Bbb{C}$. We construct a natural  ramification Cartier divisor $\mathcal{R}\subseteq \mathcal{Y}$ as follows. Let $U\to \mathcal{Y}$
be an atlas. We get a morphism of smooth schemes of the same dimension
$U\times_{\mathcal{Y}}\mathcal{X}\leto{\pi'} U$. The map on tangent spaces results in  a map of bundles of the same rank on $U\times_{\mathcal{Y}}\mathcal{X}$ given by $T (U\times_{\mathcal{Y}} \mathcal{X}) \to \pi'^* TU$. From this data we obtain  a line bundle $\det \pi'^* TU\tensor \det T (U\times_{\mathcal{Y}}\mathcal{X})^{-1}$ and a section. We denote the corresponding Cartier divisor by $\mathcal{R}$ and the line bundle by $\mathcal{O}(\mathcal{R})$. The line bundle and the section glue in the smooth topology of $\mathcal{X}$.


To prove Part (c) of Proposition \ref{bee}, we introduce a line bundle on $\widehat{\Omega}^0_S(d,r,\mn,n,\vec{I})$ with a section which we will later prove to be equal to $\mathcal{O}(\mathcal{R})$ and the corresponding section. Consider a point
$$x=(\mv,\mw,\me,\gamma)\in \widehat{\Omega}^0_S(d,r,\mn,n,\vec{I}).$$
Here,
\begin{itemize}
\item[--]$\mv\subset \mw$ is a locally free subsheaf. Let $\mathcal{Q}=\mw/\mv$ which is by assumption, locally free near points of $S$.
\item[--]Let $\mathcal{F}$ (respectively $\mathcal{G}$) be the induced flags on the fibers of $\mv$ (respectively $\mathcal{Q}$) at points $p\in S$.
\item[--] For each $p=p_j\in S$ define a subspace $T_p$ of $\Hom(\mv_p,\mq_p)$ as follows:
$$T_p=\{\phi: \mv_p\to \mq_p\mid \phi(F^p_{k})\subseteq G^p_{i_j^{k}-k}, k=1,\dots,r \}.$$
Note that $T_p$ is the tangent space of the Schubert cell ${\Omega}^0_{I^j}(E_{\bull})$ of $\Gr(r,\mw_p)$ at $\mv_p\subseteq \mw_p$ and has dimension $|\sigma_{I^j}|$.
\end{itemize}
Consider the kernel $\mk_x$ of the natural surjective mapping of sheaves on $\Bbb{P}^1$  (the right hand side is a sum of skyscraper sheaves)
\begin{equation}\label{sky}
\source{rigid-local-systems-multiplicative-eigenvalue-problem}
\shom(\mv,\mathcal{Q})\to \bigoplus_{p\in S}\frac{\Hom(\mv_p,\mq_p)}{T_p}\mid_p.
\end{equation}
The Euler characteristic $\chi(\Bbb{P}^1,\mk_x)$ is easily computed to be zero  using \eqref{timer} (note that $\mv$ is locally free), and hence we get a determinant of cohomology line (i.e., a one dimensional complex vector space) $D(\mk_x)=\det H^0(\Bbb{P}^1,\mk_x)^*\tensor  \det H^1(\Bbb{P}^1,\mk_x)$
together with a canonical section (i.e., an element in this line). This construction (as $x$ varies) gives
\begin{itemize}
 \item[-] A locally free sheaf $\mk$ on $\widehat{\Omega}^0_S(d,r,\mn,n,\vec{I})\times\Bbb{P}^1$.
\item[-] A line bundle $\Theta$ on $\widehat{\Omega}^0_S(d,r,\mn,n,\vec{I})$, together with a section $\theta$ on $\widehat{\Omega}^0_S(d,r,\mn,n,\vec{I})$. The fiber of this bundle over $x\in \widehat{\Omega}^0_S(d,r,\mn,n,\vec{I})$ is the determinant of cohomology $\mathcal{D}(\mk_x)$ of the coherent sheaf $\mk_x$ (as above) on $\Bbb{P}^1$.
\end{itemize}
\begin{proposition}
\source{rigid-local-systems-multiplicative-eigenvalue-problem}
\notready\label{torro}
\source{rigid-local-systems-multiplicative-eigenvalue-problem}
The Cartier divisor on $\widehat{\Omega}^0_S(d,r,\mn,n,\vec{I})$  corresponding to $\theta$ equals $\mathcal{R}$, the ramification divisor of the representable morphism of smooth Artin stacks $\pi:\widehat{\Omega}^0_S(d,r,\mn,n,\vec{I})\to \Par_{n,\mn,S}$.
\end{proposition}
The above equality of divisors at the level of sets is both easy (see the first paragraph below) and sufficient for our main results (the proof of Proposition \ref{bee}(c)). 
 \begin{proof}
 Note that the ramification divisor $\mathcal{R}$ is supported on the points where the fiber of $\pi$ is not smooth. The tangent space of a fiber of $\pi$ at a point $x$ is the same as $H^0(\Bbb{P}^1,\mathcal{K}_x)$, which is non zero only at points at which $\theta$ vanishes. Therefore the desired equality is true at the level of sets. We now show that we have equality as Cartier divisors.

 We ignore the data of $\det(\mw)\leto{\gamma}\mathcal{N}$ for simplicity. There exists an $N>0$, and a large open substack $\op{Bun}'_{\mn}(n)$ of $\op{Bun}_{\mn}(n)$ (complement of codimension $\geq 2$) such that
\begin{itemize}
\item[-] For points $(\mw,\gamma)\in \op{Bun}'_{\mn}(n)$, the vanishing $H^1(\Bbb{P}^1,\mw(N))=0$ holds, and $H^0(\Bbb{P}^1,\mw(N))$ is generated by global sections.
\end{itemize}
 Suppose that the $H^0(\Bbb{P}^1,\mw(N))$  above is $M$-dimensional. Let $\mt=(\mathcal{O}^{\oplus M})\tensor \mathcal{O}$. There is a suitable component $Q'$ of the quot scheme of quotients of $\mt$ such that there is a smooth atlas $Q'\to \op{Bun}'_{\mn}(n)$ so that $\op{Bun}'_{\mn}(n)$ is the stack quotient
$[Q'/G]$ with $G= \operatorname{Aut}(\mt)=\operatorname{GL}(M)$.

The space $\Par_{n,\mn,S}$ is a bundle over $\op{Bun}'_{\mn}(n)$ whose fibers are of the form $\Fl(n)^s$. We can base change  $\widehat{\Omega}^0(d,r,\mn,n,\vec{I})\to \Par_{n,\mn,S}$ to $Q'$. The base changed objects are now smooth schemes with actions of $G$, so that the original objects are stack quotients by $G$ (but only over $\op{Bun}'_{\mn}(n)$). Let the resulting map of smooth schemes be denoted by $\pi':W\to P$.

Note that $W$ parametrizes points $(\mv,\mw,\me,\gamma)\in \widehat{\Omega}^0_S(d,r,\mn,n,\vec{I})$ together with a surjection $\mt\twoheadrightarrow  \mw(N)$, so that $\mw^*$ is a subbundle of $\mt^*(N)$. Therefore $W$ is a fiber-bundle over a hyperquot scheme. $P$ is the data of $(\mw,\me,\gamma)$  with a surjection $\mt\twoheadrightarrow \mw(N)$.


We claim that given a point $x\in W$, there exist sheaves $\mathcal{A}_x$ and $\mathcal{B}_x$ on $\Bbb{P}^1$, and a surjective map of sheaves $\mathcal{A}_x\to \mathcal{B}_x$  so that
\begin{enumerate}
\item $TW_x =H^0(\Bbb{P}^1, \mathcal{A}_x)$, $TP_{\pi'(x)}=H^0(\Bbb{P}^1,\mathcal{B}_x)$ and $H^1(\Bbb{P}^1,\mathcal{A}_x)=H^1(\Bbb{P}^1,\mathcal{B}_x)=0$.
\item The sheaves $\mathcal{A}_x$ and $\mathcal{B}_x$ fit together and give sheaves $\mathcal{A}$ and $\mathcal{B}$ on $W\times \Bbb{P}^1$ together with a surjective map $\mathcal{A}\to \mathcal{B}$.
\item The sheaf $\mathcal{K}$ pulled back to $W\times\Bbb{P}^1$ equals the kernel of  $\mathcal{A}\to \mathcal{B}$. The pushforward, in the derived category, of (this pullback of) $\mk$ to $\Omega$
    is represented by a two term complex of vector bundles with fibers $H^0(\Bbb{P}^1, \mathcal{A}_x)$ and $H^0(\Bbb{P}^1,\mathcal{B}_x)$.
\end{enumerate}
This will imply that $\mathcal{D}(\mk_x)$ (together with the section) is canonically identified with $\det H^0(\ma_x)^*\otimes \det H^0(\mb_x)$ (together with the canonical section), and hence with the ramification line bundle and its section, and hence complete the proof.


To prove the claim, we appeal to the description of tangent spaces to hyperquot scheme parametrizing hyperquotients of the trivial bundles (i.e., of flags of subsheaves of the trivial bundles), see e.g., \cite{CF} (particularly Theorem 1.2, and Proposition E in the appendix). Ignoring the flag structures, both $W$ and $P$ are hyperquot schemes and the description of $\mk$ as the kernel of $\ma\to \mb$ follows from \cite[Prop. E]{CF}.

When flag structures are introduced, we need to replace $P$ by a product (locally) over $\op{Bun}'_{\mn}(n)$ of the form $\op{Bun}'_{\mn}(n)\times (\Fl(n))^s$ because the fibers at $p\in S$ of the universal bundles on $\op{Bun}'_{\mn}(n)\times \Bbb{P}^1$ are locally trivial. The replacement of $\ma_x$ will be a subsheaf $\ma'_x$ of $\ma_x\oplus T_e(\Fl(n))^s$ with the quotient a sum of skyscraper sheaf (on fibers), see Remark \ref{afternoon} below. We want $H^1$ of this sheaf to be zero to apply the above reasoning. This follows because we know that the dimension of $H^0(\ma'_x)$ is the expected one ($\widehat{\Omega}^0_S(d,r,\mn,n,\vec{I})$ is smooth of the expected dimension).
\end{proof}
\begin{remark}
\source{rigid-local-systems-multiplicative-eigenvalue-problem}
\notready\label{afternoon}
\source{rigid-local-systems-multiplicative-eigenvalue-problem}
Consider the universal Schubert variety $\Omega_I$ in $\Gr(r,n)\times\Fl(n)$, and a smooth point $x=(V,E_{\bull})$ in $\Omega_I$. Then the tangent space of $\Omega_I$ fits into an exact sequence
$$0\to T_x(\Omega_I) \to T(\Gr(r,n))_V\oplus T(\Fl(n))_{E_{\bull}}\to Q\to 0$$ where $Q$ is a vector space of dimension given by the codimension of $\Omega_I(E_{\bull})$ in $\Gr(r,n)$. The map $T(\Gr(r,n))_V\oplus T(\Fl(n))_{E_{\bull}}\to Q$ is surjective on each factor.
\end{remark}
There is a morphism  $\ParL\to {\Omega}^0(d,r,\mn,n,\vec{I})$. Let $\mathcal{R}_L\subseteq \ParL$ be the pullback of $\mathcal{R}$. From the description of $\mathcal{R}$ as the zeroes of a determinant given in \eqref{sky}, it follows that
\begin{lemma}
\source{rigid-local-systems-multiplicative-eigenvalue-problem}
\notready
The Cartier divisor $\mathcal{R}$ on  $\widehat{\Omega}^0(d,r,\mn,n,\vec{I})$ equals $\pi^*\mathcal{R}_L$ where
$\pi: \widehat{\Omega}^0(d,r,\mn,n,\vec{I})\to \ParL$.
\end{lemma}

Note that \eqref{morph1} restricts to
 \begin{equation}
\xymatrix{
 \Omega^0-\mathcal{R} \ar[r]\ar[d] &  \widehat{\Omega}^0-\mathcal{R}\ar[dl]\\
\ParL-\mathcal{R}_L
}
\end{equation}
\subsection{Proof of Proposition \ref{bee}, (c)}\label{martel}
\source{rigid-local-systems-multiplicative-eigenvalue-problem}
We need to show that a point $x\in\widehat{\Omega}_{S}^0-\widehat{\Omega}^0$, $H^0(\mk_x)\neq 0$. This is clear because $\shom(\mv,\mathcal{Q})$ has torsion supported outside of $S$, and therefore gives rise to torsion in $\mk_x$.
This proves (c) of Proposition \ref{bee}. Proposition \ref{bee} has the following important corollary:
\begin{corollary}
\source{rigid-local-systems-multiplicative-eigenvalue-problem}
\notready\label{carole}
\source{rigid-local-systems-multiplicative-eigenvalue-problem}
Suppose $\langle\sigma_{I^1}, \sigma_{I^2},\dots,\sigma_{I^s}\rangle_{d,D}=1$, and $1\not\in I^i$ for all $i$. Then $\widehat{\Omega}^0(d,r,\mn,n,\vec{I})-\mathcal{R}$ maps isomorphically to an open substack of $\Par_{n,\mn,S}$ whose complement has codimension $\geq 2$.
\end{corollary}
\begin{proof}
 The assertion follows from Propositions \ref{maya} and \ref{bee}, the surjectivity and birationality of $\Omega\to \Par_{n,\mn,S}$, and Zariski's main theorem.
\end{proof}
